\problem{Стабильность звезды}{10.0}

\pic{9}{R}{0pt}{0.263}{figures/2.1.png}

Звезда --- не просто одна из тысячи ярких точек на ночном небосводе, но ещё и важный объект для астрофизических исследований. Эта задача нацелена изучить базовые механизмы образования и сохранения равновесия звезды. Гравитационная постоянная $G=\qty{6.674e-11}{\Ggrav}$, скорость света $c=\qty{3.00e8}{\m\per\s}$, постоянная Больцмана $k=\qty{1.38e-23}{\J\per\K}$, постоянная Стефана--Больцмана $\sigma=\qty{5.67e-8}{\W\per\m\squared\per\K\tothe4}$, атомная единица массы $\qty{1}{\amu}=\qty{1.66e-27}{\kg}$.

В астрофизике индекс ``$\odot$'' говорит о том, что данная величина относится к нашему Солнцу. Для данной задачи считайте известными: масса Солнца $M_\odot=\qty{2e30}{\m}$, радиус Солнца $R_\odot=\qty{6.96e8}{\m}$, а также средний орбитальный радиус Земли $R_E=\qty{1.5e11}{\m}$.

\part{Параллакс}

\pic{10}{R}{0pt}{0.25}{figures/2.2.png}

Звёздный параллакс --- самый первый метод по определению расстояния звезды от нашей Солнечной системы. Так, например, спутник Hipparcos (\textbf{Hi}gh \textbf{P}recision \textbf{Par}allax \textbf{Co}llecting \textbf{S}atellite), запущенный в 1989 году для сбора астрометрических данных, может обнаруживать параллаксы с точностью до миллиарксекунд. Рассмотрим звезду $S$, которая находится перпендикулярно орбитальной плоскости Земли, как показано на рисунке справа (такой, например, является $\delta$ Дракона). При наблюдении за ней в течение одного года она ``движется'' по кругу $C$ относительно статичных дальних звёзд. Угловой диаметр $\Delta\varphi$ орбиты $C$, согласно Hipparcos, равняется $\num{66.7e-3}$ угловым секундам.

\begin{task}
    \item Определите расстояние $d$ от Солнца до звезды $S$ в световых годах.
\end{task}

\part{Формирование протозвезды}

\pic{10}{R}{0pt}{0.37}{figures/2.3.png}[Туманность W51 --- одна из самых больших фабрик звёзд в Млечном пути]

Большинство звёзд образуются кластерами, среди очень холодных (порядка 10--30 К) молекулярных облаков в межзвёздном пространстве. Под действием гравитационного притяжения, газовое облако начинает сжиматься, формируя протозвезду. Рассмотрим упрощённую модель водородного облака в виде однородного шара радиуса $R$ (который можно считать известным в пунктах \taskref{2}--\taskref{4}), а также при средней температуре $T=\qty{10}{\K}$ (при таких температурах вращательные степени свободы молекул ``заморожены''). Средняя масса каждой молекулы шара равна $m=\qty{2}{\amu}$, их концентрация $n=\qty{4e8}{\m\tothe{-3}}$. В действительности, механизм формирования протозвезды гораздо более сложный: коллапсу сопротивляются давление газа, наличие магнитного поля, центробежные силы и турбулентность, и потому коллапс обычно запускается распространением ударных волн.

\begin{task}
    \item Выразите полную внутреннюю энергию $U$ теплового движения молекул облака.
    \item Определите, как ускорение свободного падения $g(r)$ зависит от радиуса $r$ внутри облака и полной его массы $M$.
    \item Собственная гравитационная энергия $E_G$ шара массы $M$ и радиуса $R$ имеет вид
    $$
    E_G = -f\frac{GM^2}{R},
    $$
    где $f$ --- безразмерный коэффициент, в общем случае показывающий распределение массы внутри шара. Определите $f$ в случае однородного шара.
    \item Обязательным условием коллапса облака является отрицательность полной механической энергии $E$. Отсюда определите критическую полную массу $M_J$ облака при данных числовых параметрах в единицах $M_\odot$. Эта величина называется массой Джинса --- оттуда и индекс $J$.
\end{task}

\part{Параметры Солнца}

Сфокусируем наше внимание на Солнце. Как и прочие звёзды, оно не коллапсирует под собственным гравитационным притяжением ввиду наличия градиента давлений $P(r)$ внутри себя. В данной части пренебрегите вкладом радиации Солнца в $P(r)$, а также конвективными эффектами. \textbf{Внимание:} результаты пунктов \taskref{2}--\taskref{5} недействительны далее, так как распределение вещества внутри Солнца \textbf{нельзя} считать однородным!

\begin{task}
    \item Рассмотрите гидростатическое равновесие элемента Солнца в виде диска плотности $\rho(r)$ и толщины $\Delta r$ на расстоянии $r<R_\odot$. Оттуда выведите зависимость радиального градиента давления $dP/dr$ как функции от $r$, $\rho(r)$ и части массы Солнца $M(r)$, заключённой внутри сферической оболочки радиуса $r$ (таким образом, $M(R_\odot)=M_\odot$).
\end{task}

Поскольку давление и плотность являются сложными функциями от расстояния, можно упростить расчёты, если использовать $\langle P\rangle$ --- давление Солнца, усреднённое по его объёму:
$$
\langle P\rangle = \frac1{V_\odot} \int_0^{V_{\odot}} P\, dV.
$$
Здесь $V_\odot$ --- полный объём Солнца. \textit{Подсказка:} в пункте \taskref{7} Вам может пригодиться метод интегрирования по частям $\int_a^b u\,dv= uv|_a^b-\int_a^b v\,du$.

\begin{task}
    \item Покажите, что полная гравитационная энергия $E_G$ Солнца и среднее давление $\langle P\rangle$ связаны выражением
    $$
    E_G = \beta\langle P\rangle V_\odot,
    $$
    где $\beta$ --- безразмерный численный параметр. Найдите значение $\beta$.
    \item Ввиду того, что плотность Солнца растёт по мере приближения к центру, коэффициент $f$ к выражению для $E_G$ отличается от ответа в \taskref{4} для однородного облака. В предположении, что $f\approx1$ для Солнца, оцените численное значение среднего давления $\langle P\rangle$ внутри Солнца.
\end{task}

Оценим среднюю температуру $T_0$ внутри Солнца. Для этого допустим, что давление $\langle P\rangle$ и плотность распределены однородно внутри Солнца, и что вещество Солнца является идеальным газом из электронов и ионизованных водорода с гелием. Средняя масса частиц ионизованного газа примерно равна $\overline m=\qty{0.61}{\amu}$.

\begin{task}
    \item С учётом допущений выше оцените численно температуру $T_0$ внутри Солнца.
\end{task}

\part{Излучение Солнца}

В финальной части этой задачи изучим излучательные свойства Солнца. Известно, что Солнце, как и большинство прочих звёзд, имеет спектр теплового излучения очень близкий к спектру абсолютного чёрного тела, а светимость (т.е. полная излучательная мощность) Солнца равна $L_\odot=\qty{4e26}{\W}$.
\begin{task}
    \item Оцените температуру $T_s$ на поверхности Солнца.
\end{task}

\pic{11}{R}{0pt}{0.27}{figures/2.4.png}

В первом приближении можно рассмотреть Солнце как шар из электронов с ионами, которые находятся в термодинамическом равновесии с радиацией при температуре $T_0$ --- тогда бы Солнце было абсолютно чёрным радиатором при $T_0$ и, соответственно, другой светимости $L_\odot'$. Различие между $L_\odot$ и $L'_\odot$ появляется ввиду того, что радиация внутри Солнца постоянно рассеивается, поглощается и излучается электронами и ионами.

Чтобы смоделировать явление \textit{радиативной диффузии}, представим, что фотон, начавший движение в центре Солнца, идёт по ломаной линии ввиду поглощения/излучения, причём среднюю длину свободного пробега $l$ можно считать постоянной для фотона. Направление следующего смещения $\vec l_i$ перед следующим поглощением совершенно случайно и не зависит от предыдущего перемещения $\vec l_{i-1}$, и при этом $|\vec l_i|=l$ для всех $i$.

\begin{task}
    \item Запишите полный вектор перемещения $\vec D$ фотона через $\vec l_1,\dots,\vec l_N$. Чему равно среднее перемещение $\langle\vec D\rangle$ спустя $N$ шагов смещения фотона?
    \item Оцените среднеквадратичное смещение $\langle D^2\rangle$ в зависимости от $N$ и $l$. Отсюда найдите оценочное время $t$ для того, чтобы фотон покинул Солнце.
\end{task}

Ввиду радиативной диффузии мы получаем, что $L_\odot \approx L_\odot'\cdot l/R_\odot$. Это соотношение поможет нам понять, как светимость $L_\odot$ зависит от массы Солнца $M_\odot$.

\begin{task}
    \item Оцените численно значение $l$ в зависимости от $T_0$, $T_s$, и $R_\odot$. Отсюда оцените значение $t$ из пункта \taskref{12}. Результат подтверждает, что Солнце внутри весьма непрозрачно.
    \item Используя результат пункта \taskref{9} покажите, что $L_\odot$ и $M_\odot$ связаны между собой как
    $$
    L_\odot = C\cdot\langle\rho\rangle\cdot(M_\odot)^n.
    $$
    Здесь $\langle\rho\rangle=\frac{M_\odot}{V_\odot}$. Покажите, без численных расчётов, связь коэффициента $C$ с параметрами $\overline m$, $l$, а также фундаментальными константами, и определите численно значение степени $n$.
\end{task}

Интересно отметить важность результата \taskref{14}. Степенная зависимость $L_\odot$ от $M_\odot$ достаточно хорошо наблюдается для звёзд главной последовательности. График зависимости масса--светимость для звёзд в интервале $0.2M_\odot$--$50M_\odot$ показан ниже ($\log$ берётся по основанию 10).

\begin{task}
    \item На основании графика оцените степень $n$, определённый как в пункте \taskref{14}, без оценки погрешности.
\end{task}

\cpic{1}{figures/2.5.png}[Зависимость масса--светимость в интервале $0.2M_\odot$--$50M_\odot$. (Torres et al. (2010))]
